\begin{center}\begin{tikzpicture}[font=\small,>=stealth,line join=round]
\draw[gray!40,rounded corners=3pt] (0,0) rectangle (14.2,8.3);\node at (7.1,7.700000000000001) {外层没有奇偶性，内偶仍能保证整体偶};
\begin{scope}[shift={(4.1,2.068181818181818)},x=1.875cm,y=0.8363636363636363cm]\draw[->,gray] (-1.6,0)--(1.6,0) node[right]{$x$};\draw[->,gray] (0,-0.5)--(0,5) node[above]{$y$};\draw[blue,thick,solid,domain=-1.4:1.4,samples=180] plot(\x,{pow(2,(\x)^2)});\fill[orange!85!black] (-1,2) circle(2.2pt);\fill[orange!85!black] (1,2) circle(2.2pt);\end{scope}\node[anchor=west,align=left,text width=5.0cm] at (8.4,4.0) {$F=2^{x^2}$\\[10pt]$x$ 与 $-x$ 先得到同一指数\\[10pt]外层 $2^u$ 非奇非偶\\[10pt]最后函数值仍然相同};
\node[text width=12.6cm,align=center] at (7.1,.65) {两路先合并为同一输入，才是“内偶则偶”的决定因素。};
\end{tikzpicture}\end{center}
